\section{Discussion}
\label{sec:discussion}

Our results suggest a nuanced conclusion. On the one hand, raw last-layer attention is not a reliable standalone proxy for token importance. Its agreement with both gradient $\times$ input and LOO is weak, and its deletion behavior is consistently less disruptive than LOO. If one interprets token importance as ``which token most changes the prediction when removed,'' then attention is clearly not the best choice in this setting.

On the other hand, attention is not completely detached from model behavior. Attention-based deletion is stronger than random deletion throughout the fixed top-$k$ experiments, which means it captures some non-random signal about influential tokens. In our setup it also outperforms gradient $\times$ input on both confidence drop and flip rate. This matters because it prevents an overly strong conclusion. The empirical result is not that attention is meaningless; rather, it is that attention is incomplete and unreliable when used as a direct substitute for stronger importance measures.

One plausible explanation is that raw Transformer attention mixes several functions at once. Prior analyses of BERT attention show that heads often focus on delimiters, punctuation, or positional patterns in addition to semantically meaningful content~\cite{clark2019bert,sun2021effective}. Averaging last-layer attention from \texttt{[CLS]} may therefore preserve some useful task signal while also inheriting structural attention behaviors that do not correspond closely to causal influence on the output.

The project also has several limitations. First, it uses only one task and one model family, so the conclusions should not be generalized automatically to other datasets or larger language models. Second, the analysis is performed at the wordpiece level, which can split meaningful words into multiple ranked tokens. Third, we examine only one attention summary: mean last-layer \texttt{[CLS]} attention. Different layers, individual heads, or effective attention variants may behave differently. Finally, LOO itself is a local perturbation method rather than a complete causal account, and it is computationally more expensive than attention or gradient-based scoring.
